\documentclass[11pt]{article}
\usepackage[letterpaper,left=18mm,right=18mm,top=20mm,bottom=18mm,headheight=14pt]{geometry}
\usepackage{tikz,xcolor,fancyhdr,lmodern}
\usetikzlibrary{arrows.meta,decorations.pathreplacing}
\renewcommand{\familydefault}{\sfdefault}
\pagestyle{fancy}\fancyhf{}\renewcommand{\headrulewidth}{0.4pt}
\newcommand{\bandname}{Grades 2--3}
\fancyhead[L]{\small Week 92 / Fold-and-stretch predictions / \bandname}
\fancyfoot[L]{\scriptsize Bellingham Math Circle / Week 92 / W92-S-v2}
\fancyfoot[R]{\scriptsize\thepage}
\setlength{\parindent}{0pt}\setlength{\parskip}{0pt}
\pdfinfo{/Title (Week 92 - Fold-and-stretch predictions) /Author (Bellingham Math Circle) /CreationDate (D:20261017000000Z) /ModDate (D:20261017000000Z)}
\pdftrailerid{<92A113BA84D2123FF010701001010101><92A113BA84D2123FF010701001010101>}
\definecolor{ink}{HTML}{203642}\definecolor{bandfill}{HTML}{B8D9E6}
\tikzset{every node/.style={font=\sffamily\small,text=ink},arr/.style={-{Stealth[length=2mm]},thick,draw=ink}}
\newcommand{\txt}[3]{\node[anchor=north west,align=left,text width=174mm,inner sep=0] at (0,#1) {#2#3};}
\newcommand{\prob}[3]{\txt{#1}{\textbf{Problem #2:} }{#3}}
\newcommand{\track}[2]{\begin{scope}[shift={(8,#1)}]\draw[ink] (0,0)--(160,0);\foreach \i in {0,...,16}{\draw[ink] (10*\i,-1.5)--(10*\i,1.5);\node[font=\scriptsize,anchor=north] at (10*\i,-2){\i};}\node[anchor=south west,font=\scriptsize] at (0,3){#2};\end{scope}}
\newcommand{\range}[3]{\fill[fill=bandfill,fill opacity=0.55] (8+10*#2,#1-1.3) rectangle (8+10*#3,#1+1.3);\fill[ink] (8+10*#2,#1) circle(0.7);\fill[ink] (8+10*#3,#1) circle(0.7);}
\newcommand{\point}[2]{\fill[ink] (8+10*#2,#1) circle(1.25);}
\newcommand{\canvas}{\begin{tikzpicture}[x=1mm,y=1mm]\path[use as bounding box] (0,-231) rectangle (178,0);}
\newcommand{\finishcanvas}{\end{tikzpicture}}

% Reusable continuous mapping aid: horizontal input, right distance, left output.
% The paired vertical scales encode the complete doubling operation.
\newcommand{\movemap}[2]{
\begin{scope}[shift={(12,#1)}]
\draw[ink] (0,0)--(160,0);
\foreach \i in {0,...,16}{\draw[ink] (10*\i,-1)--(10*\i,1);\node[font=\scriptsize,anchor=north] at (10*\i,-2){\i};}
\foreach \i in {0,...,8}{
\draw[gray!20] (0,2*#2*\i)--(160,2*#2*\i);
\draw[ink] (-1,2*#2*\i)--(1,2*#2*\i);
\node[font=\scriptsize,anchor=east] at (-2,2*#2*\i){\pgfmathparse{int(2*\i)}\pgfmathresult};
\draw[ink] (159,2*#2*\i)--(161,2*#2*\i);
\node[font=\scriptsize,anchor=west] at (162,2*#2*\i){\i};}
\draw[ink] (0,0)--(0,16*#2);\draw[ink] (160,0)--(160,16*#2);
\draw[ink,very thick] (0,0)--(80,16*#2)--(160,0);
\node[anchor=south west,font=\scriptsize] at (-10,16*#2+3){Output};
\node[anchor=south east,font=\scriptsize] at (165,16*#2+3){Nearer-end distance};
\node[anchor=north,font=\scriptsize] at (80,-8){Input};
\end{scope}}
\begin{document}
\canvas
\txt{0}{}{One move: find the nearer-end distance in track units, then double it. On the map, go up from your input to the sloping line, then across to read the distance and output. Put the marker at the output on the 0--16 track. Points between ticks are allowed.}
\movemap{-123}{5}
\draw[arr,dashed] (142,-123)--(142,-93);
\draw[arr,dashed] (142,-93)--(172,-93);
\draw[arr,dashed] (142,-93)--(12,-93);
\fill[ink] (142,-123) circle(1);\fill[ink] (142,-93) circle(1);
\node[fill=white,inner sep=1.5mm,font=\scriptsize] at (89,-110){input 13 $\longrightarrow$ distance 3 $\longrightarrow$ output 6};\prob{-141}{1}{Put two markers at different whole-number starts. Find every pair that meets after one move. Then find two starts whose markers move farther apart for three moves in a row.}
\track{-177}{ }\track{-206}{ }\finishcanvas
\newpage
\canvas
\txt{0}{}{A band includes every position between its ends, including the ends. Move all of them through the map. One marker follows the rule from its unknown start; it is never placed at a fresh start.}
\movemap{-83}{3}
\fill[bandfill] (62,-84.3) rectangle (112,-81.7);
\draw[ink,line width=1.3mm] (62,-53)--(92,-35)--(112,-47);
\draw[ink,line width=1.3mm] (12,-53)--(12,-35);
\draw[ink,line width=1.3mm] (172,-53)--(172,-35);
\foreach \x/\y in {62/-53,92/-35,112/-47}{\draw[ink,dashed] (\x,-83)--(\x,\y);\draw[ink,dashed] (12,\y)--(172,\y);\fill[ink] (\x,-83) circle(0.8);\fill[ink] (\x,\y) circle(0.8);}
\node[anchor=north,font=\scriptsize] at (90,-98){Input band 5--10 $\longrightarrow$ distance band 5--8 $\longrightarrow$ output band 10--16};\prob{-110}{2}{Shade all possible places after one move for each starting band.}
\track{-131}{Start}\range{-131}{1}{4}\track{-144}{After one move}
\track{-160}{Start}\range{-160}{10}{15}\track{-173}{After one move}
\track{-189}{Start}\range{-189}{6}{11}\track{-202}{After one move}
\track{-218}{Start}\range{-218}{2}{14}\track{-231}{After one move}\finishcanvas
\newpage
\canvas
\prob{0}{3}{Start with each band below. Move the whole band until every place on the track is possible. How many moves does each need?}
\track{-32}{Start}\range{-32}{1}{2}
\track{-51}{After 1 move}\track{-70}{After 2 moves}\track{-89}{After 3 moves}\track{-108}{After 4 moves}
\track{-139}{Start}\range{-139}{7}{9}
\track{-158}{After 1 move}\track{-177}{After 2 moves}\track{-196}{After 3 moves}\track{-215}{After 4 moves}\finishcanvas
\newpage
\canvas
\gdef\bandname{Grades 4--5}
\txt{0}{}{Use the move and band rules on pages 1--2.}
\prob{-13}{4}{A child moves just the ends of the band from 6 to 10. Both ends land at 12. Is 12 the only possible next place? Shade the whole output band. Find two more starting bands for which shading only between the two endpoint outputs misses possible places.}
\track{-51}{Start}\range{-51}{6}{10}\track{-70}{After one move}
\track{-94}{Start}\track{-113}{After one move}
\track{-137}{Start}\track{-156}{After one move}\prob{-175}{5}{Find a starting band that becomes wider after one move and a band that keeps the same width. Could any band become narrower? Explain with a drawing.}
\draw[gray!35] (8,-223)--(168,-223);\finishcanvas
\newpage
\canvas
\prob{0}{6}{A marker starts somewhere in a band with one end at 0. Make the widest starting band you can that keeps every possible position between 0 and 2 through the first 3 moves. Then try 4 moves and 6 moves.}
\track{-34}{The whole track}
\draw[ink,very thick] (8,-34)--(18,-34);
\fill[ink] (13,-34) circle(0.8);
\draw[gray,dashed] (8,-39)--(8,-62);\draw[gray,dashed] (18,-39)--(168,-62);
\draw[ink] (8,-62)--(168,-62);
\foreach \x/\lab in {8/0,88/{$\frac12$},168/1}{\draw[ink] (\x,-64)--(\x,-60);\node[anchor=north,font=\scriptsize] at (\x,-65){\lab};}
\fill[ink] (88,-62) circle(1);
\draw[ink,dashed] (13,-39)--(88,-60);
\node[anchor=west,fill=white,inner sep=1mm,font=\scriptsize] at (67,-48){The same position: $\frac12$};\txt{-77}{}{Mark starting bands on these enlarged 0--1 lines. Moves still use the original 0--16 rule.}
\foreach \y/\lab in {-104/3 moves,-139/4 moves,-174/6 moves}{
\node[anchor=south west] at (8,\y+3){\lab};
\draw[ink] (8,\y)--(168,\y);
\draw[ink] (8,\y-2)--(8,\y+2);\draw[ink] (168,\y-2)--(168,\y+2);
\node[anchor=north,font=\scriptsize] at (8,\y-3){0};\node[anchor=north,font=\scriptsize] at (168,\y-3){1};}\prob{-193}{7}{Someone names a number of moves. Can you always choose a starting band with some width that keeps every possible position between 0 and 2 for that long? Explain.}
\finishcanvas
\newpage
\canvas
\prob{0}{8}{Move each labelled piece separately. Shade its possible places after the stated number of moves on the matching track. Does the pattern continue if you keep cutting every piece in half? Explain.}
\foreach \y/\n/\lab/\first in {-36/2/1 move/65,-65/4/2 moves/67,-94/8/3 moves/71}{
\node[anchor=south west,font=\scriptsize] at (8,\y+10){\lab};
\foreach \k in {1,...,\n}{\pgfmathsetmacro{\xa}{8+160*(\k-1)/\n}\pgfmathsetmacro{\xb}{8+160*\k/\n}
\fill[bandfill] (\xa,\y-1.3) rectangle (\xb,\y+1.3);
\fill[ink] (\xa,\y) circle(0.7);\fill[ink] (\xb,\y) circle(0.7);
\pgfmathtruncatemacro{\letter}{\first+\k-1}
\node[font=\scriptsize,anchor=south] at ({(\xa+\xb)/2},\y+2){\char\letter};}
\track{\y}{ }}
\foreach \col/\first in {0/65,1/72}{
\foreach \r in {0,...,6}{
\pgfmathtruncatemacro{\letter}{\first+\r}
\begin{scope}[shift={({8+88*\col},{-120-14*\r})}]
\node[anchor=east,font=\scriptsize] at (-2,0){\char\letter};
\draw[ink] (0,0)--(70,0);
\foreach \i in {0,...,16}{\draw[ink] (70*\i/16,-1)--(70*\i/16,1);}
\foreach \i in {0,4,8,12,16}{\node[anchor=north,font=\tiny] at (70*\i/16,-2){\i};}
\end{scope}}}\finishcanvas
\newpage
\canvas
\prob{0}{9}{Will every starting band with some width eventually make the whole track possible? Include bands whose ends lie between ticks in your explanation.}
\track{-40}{ }\track{-70}{ }\track{-100}{ }\finishcanvas
\end{document}
